\section{Постановка задачи}

\begin{frame}[c]{Три парадигмы обучения}
  \begin{columns}[T]
    \begin{column}{0.31\textwidth}
      \begin{block}{С учителем}
        Для каждого объекта известен правильный ответ. Учимся предсказывать ответ для новых объектов.

        \smallskip
        \lead{Пример:} диагноз по симптомам.
      \end{block}
    \end{column}
    \begin{column}{0.31\textwidth}
      \begin{block}{Без учителя}
        Есть только данные, ответов нет. Ищем структуру: например, группы похожих объектов.

        \smallskip
        \lead{Пример:} группировка новостей по сюжетам.
      \end{block}
    \end{column}
    \begin{column}{0.31\textwidth}
      \begin{block}{С подкреплением}
        Агент действует в среде и получает награду за удачные действия.

        \smallskip
        \lead{Пример:} робот учится запрыгивать на стол.
      \end{block}
    \end{column}
  \end{columns}
\end{frame}

\begin{frame}[c]{Что предсказывает модель?}
  \lead{Вопрос на подумать.}
  \begin{numrows}
        1 & Время поездки на каршеринге\\
        2 & Кликнет ли пользователь по рекламе\\
        3 & Тема научной статьи: математика, биология, психология\ldots\\
        4 & Теги ресторана: «кофе», «веранда», «живая музыка»\\
        5 & Порядок документов в поисковой выдаче\\
      \end{numrows}\par
  \medskip
  \textbf{Чем различаются ответы? Какие группы задач вы видите?}
\end{frame}

\begin{frame}[c]{Обучение с учителем: виды задач}
  \small
  \zebra{2}\small
  \begin{tabular}{>{\raggedright\arraybackslash}p{0.286\textwidth}>{\raggedright\arraybackslash}p{0.593\textwidth}}
    \rowcolor{acDark} \textcolor{white}{\textbf{Задача}} & \textcolor{white}{\textbf{Что предсказываем}}\\
    \textcolor{acMid}{\textbf{Регрессия}} & Таргет --- число или вектор. Время поездки на каршеринге, спрос на товар.\\
    \textcolor{acMid}{\textbf{Бинарная классификация}} & Да или нет. Кликнет ли по рекламе, вернёт ли кредит.\\
    \textcolor{acMid}{\textbf{Многоклассовая}} & Один класс из нескольких. Тема научной статьи.\\
    \textcolor{acMid}{\textbf{Многометочная}} & Несколько классов сразу, они могут пересекаться. Теги ресторана.\\
    \textcolor{acMid}{\textbf{Ранжирование}} & Упорядочить объекты. Поисковая выдача по релевантности.\\
    \textcolor{acMid}{\textbf{Сложные ответы}} & Сегментация изображения, перевод текста, генерация новых объектов.\\
  \end{tabular}
\end{frame}

\begin{frame}[c]{Определите тип задачи}
  \lead{Вопрос на подумать.}
  \begin{numrows}
        1 & Предсказать курс евро к доллару на следующий день\\
        2 & Стилизовать текст: «Пиппина и Мерри похитили!» $\to$ «Граждане Тук и Брендибак были похищены неустановленными лицами»\\
        3 & Определить, котик или собачка на картинке\\
        4 & Научить робокота запрыгивать на стол из любой позы\\
        5 & Найти наборы товаров, которые покупают вместе\\
      \end{numrows}\par
\end{frame}

\begin{frame}[c]{Ответы: тип задачи}
  \small
  \zebra{2}
  \begin{tabular}{>{\raggedright\arraybackslash}p{0.267\textwidth}>{\raggedright\arraybackslash}p{0.611\textwidth}}
    \rowcolor{acDark} \textcolor{white}{\textbf{Задача}} & \textcolor{white}{\textbf{Тип}}\\
    1. Курс евро & Регрессия: предсказываем число\\
    2. Стилизация текста & Генерация новых объектов на основе существующих\\
    3. Котик или собачка & Бинарная классификация: два класса\\
    4. Робокот & Обучение с подкреплением: перебирать траектории и узнавать ответ слишком долго и дорого\\
    5. Наборы товаров & Обучение без учителя: ответов нет\\
  \end{tabular}
\end{frame}
